\documentclass[conference]{IEEEtran}

% ---- Packages -------------------------------------------------------
\usepackage[T1]{fontenc}
\usepackage[utf8]{inputenc}
\usepackage{cite}
\usepackage{amsmath,amssymb,amsfonts}
\usepackage{algorithmic}
\usepackage{graphicx}
\usepackage{textcomp}
\usepackage{xcolor}
\usepackage{booktabs}       % professional tables
\usepackage{multirow}
\usepackage{url}
\usepackage{microtype}
\usepackage{subcaption}
\usepackage{hyperref}
\hypersetup{hidelinks}

% ---- Document -------------------------------------------------------
\begin{document}

\title{Beyond Mesonet4:\\
Benchmarking Modern Deepfake Detection Across Generational Gaps}

\author{
  \IEEEauthorblockN{Ryan Powers}
  \IEEEauthorblockA{\textit{Bellini College of Artificial Intelligence, Cybersecurity, and Computing}\\
  University of South Florida\\
  rp194@usf.edu}
}

\maketitle

% =====================================================================
\section{Project Goal}
\label{sec:goal}

The goal of this project was to explore how MesoNet4 — one of the earliest deepfake detectors
— compare to modern deep learning architectures and techniques
when tasked with detecting fakes across both GAN-based and
diffusion-based generation methods. Deepfake and manipulated media pose a serious threat to both national and global security. With modern technology being democratized and more accessible than ever, familiarizing ourselves and experimenting with technologies to counter media manipulation to preserve authenticity and integrity of information is becoming increasingly important. 

To begin our research into this topic, we implemented Mesonet4 as our baseline architecture. After understanding the capabilities, and limitations, of this legacy model for deepfake detection, we continued our research by experimenting with alternative CNN-based architectures and features, and than used this to compare against with modern day SOTA designs such as ViT-b16. The ultimate goal of this project was to explore the traditional techniques used to detect manipulated media and the improvements made over the years, and use this knowledge to continue our research in the same trajectory as the constant advancements being used for this purpose. Additionally, we sought to find the architectures or features that \emph{generalises across multiple generational forgery methods}. We seek to determine:
\begin{enumerate}
  \item Can a traditional CNN detector (Mesonet4/MesoInception4) remain effective against modern-day deepfake techniques?
  \item Does the addition of modern features or architectures improve detection and generalization across forgery methods?
  \item Will the addition of modern features or techniques to legacy CNN-based architectures result in competitive performance against modern architectures?
\end{enumerate}

% =====================================================================
\section{Background: Deepfake Generations and Techniques}
\label{sec:background}

Deepfake technology has evolved across three distinct generations, each exploiting a different synthesis mechanism and leaving a characteristically different forensic artifact class. Understanding these generational differences is essential context for interpreting the generalisation failures documented in Section~\ref{sec:comparison}.

\subsection{Generation 1 --- Encoder-Decoder Face Swaps (2017--2019)}

The first wave of deepfakes emerged from shared encoder / per-identity decoder autoencoders. Two autoencoders share a common encoder but train separate decoders, one per identity. At inference, a source face is encoded into the shared latent space and reconstructed by the target's decoder, producing an identity swap that is then composited onto the original frame.

FaceForensics++ \cite{rossler2019faceforensics} captures this era through four manipulation methods:

\begin{itemize}
  \item \textbf{Deepfakes}: Identity-swap via shared encoder and per-identity decoder. Artifacts manifest as blending seams at the face--background boundary, lighting discontinuities, and warping distortions around the jaw and hairline where the synthetic face region is composited.

  \item \textbf{FaceSwap}: Replaces the face region via 3D mesh reconstruction and texture transfer. Artifacts include resolution mismatches at the splice boundary and skin tone inconsistencies.

  \item \textbf{Face2Face}: Transfers facial expressions from a source actor to a target using a 3D morphable model (3DMM) for tracking and retargeting, without swapping identity. Artifacts appear as texture stitching inconsistencies at the mouth and eye regions.

  \item \textbf{NeuralTextures}: Learns a neural texture map conditioned on a 3D face model and renders it warped to the target pose. This produces the most subtle artifacts of the four methods, localised primarily to the mouth region, which is why NeuralTextures is the hardest method to detect across all architectures evaluated (Table~\ref{tab:permethod_ff}).
\end{itemize}

The unifying forensic signature of Generation 1 is \emph{spatial locality}: artifacts concentrate at blending boundaries and warped regions. CNNs with local receptive fields are inherently well-suited to detect these features.

\subsection{Generation 2 --- Higher-Fidelity Commercial GANs (2019--2021)}

Commercial deepfake pipelines improved over Generation 1 by adding post-processing stages --- colour correction, sharpening, and compression-aware blending --- that remove the most obvious blending seams. Celeb-DF-v2 \cite{li2020celeb} is representative of this generation. The forgeries in Celeb-DF-v2 are substantially more realistic than FF++ manipulations, which explains why every model in this study suffers a significant AUC drop when evaluated on it despite both datasets being GAN-based. This cross-generation failure within the GAN era foreshadows the far more severe breakdown observed against diffusion fakes.

\subsection{Generation 3 --- Diffusion-Based Synthesis (2021--Present)}

Diffusion models represent a fundamental paradigm shift. Rather than encoding and compositing a face region onto an existing frame, these models iteratively denoise a Gaussian noise signal conditioned on text prompts, reference images, or face embeddings. The synthetic face is generated \emph{globally} rather than spliced, eliminating the blending boundary that is the primary detection cue for Generation 1 methods.

This wholesale change in synthesis mechanism produces an entirely different artifact class:

\begin{itemize}
  \item \textbf{No blending boundaries.} The entire face (or image) is synthesised from noise. There is no compositing seam for local CNN filters to detect.

  \item \textbf{Over-regularised texture.} Diffusion priors encourage smooth, perceptually natural surfaces. Skin pores, fine hair, and high-frequency micro-textures are frequently over-smoothed or replaced with statistically plausible but physically inconsistent patterns.

  \item \textbf{Spectral fingerprints.} Prior work has shown that generative models leave systematic signatures in the frequency domain --- periodic patterns in the DCT/FFT spectrum arising from convolutional up-sampling in the generator or denoising U-Net backbone \cite{durall2020watch,frank2020leveraging}. These signatures differ between GAN-based and diffusion-based synthesis, though they are most visible in uncompressed or lightly compressed images and may be obscured by video codec compression (as in FF++).

  \item \textbf{Missing sensor noise.} Real camera images contain ISO grain, lens distortion, and sensor-specific imaging fingerprints. Diffusion-generated images, trained on internet-scraped data, frequently lack consistent imaging-system signatures.
\end{itemize}

Diffusion techniques evaluated in this work include Stable Diffusion 1.5 (text-to-image and inpainting), SDXL, SDXL-Refine, DreamBooth, Imagic, HPS, and CoDiff (Section~\ref{sec:dataset}). That every model trained on GAN-era FF++ collapses to near-chance AUC on diffusion fakes (Section~\ref{sec:comparison}) directly demonstrates that the spatially localised artifacts learned from Generation 1 methods are simply absent in Generation 3 outputs.

\subsection{Modern Deepfakes (2024 and Beyond)}

The current frontier extends beyond still images into real-time, multi-modal, and video-native synthesis:

\begin{itemize}
  \item \textbf{Video diffusion models} (e.g., Sora, Runway Gen-3, Kling) generate temporally consistent video sequences from text or image conditioning, eliminating the per-frame flickering that was historically a strong detection cue.

  \item \textbf{Audio-visual synthesis} (e.g., HeyGen, D-ID, ElevenLabs) combines face reenactment with neural voice cloning to produce fully synthetic video calls and interviews. Spatial-only detectors are irrelevant in this regime; detection requires cross-modal consistency analysis.

  \item \textbf{Real-time face swapping} (e.g., DeepFaceLive) operates at interactive frame rates for live video impersonation. Artifacts in this regime are dominated by temporal aliasing and latency-induced jitter rather than spatial blending errors.

  \item \textbf{Identity fine-tuning} (e.g., DreamBooth, LoRA fine-tunes of SDXL) generates photorealistic likenesses of specific individuals from fewer than ten reference photographs. Output is near-indistinguishable from real and leaves no consistent artifact pattern for a detector to exploit.
\end{itemize}

The generalisation failures documented in this work directly foreshadow the challenge of modern deepfake detection: a detector trained on any prior generation is likely to fail silently on the next, because each generation exploits a different synthesis pathway and leaves a different artifact signature.

% =====================================================================
\section{Input and Ground Truth}
\label{sec:input}

\subsection{Input Format}
All images are resized to $256 \times 256$ pixels (3-channel RGB).
This resolution is dictated by the MesoNet architecture \cite{afchar2018mesonet},
which was designed for face-crop inputs at this scale.  Images whose
shortest side is smaller than 256\,px are \emph{excluded} rather than
upscaled, as upscaling introduces artificial high-frequency artifacts
that can bias CNN feature maps. The only change from our proposal was that for our foundation model, ViT-B16 \cite{dosovitskiy2020vit}, required the images to be resized to 224px, which was done inline with fine-tuning and evaluation.

\subsection{Ground Truth Labels}
Binary labels are used throughout:
\begin{itemize}
  \item \textbf{0 — Real}: unmanipulated face images from real video
        frames (FF++) or real photographs (IMDB-WIKI via
        DeepFakeFace, Celebdfv2 Youtube real).
  \item \textbf{1 — Fake}: synthetically generated or manipulated face
        images from any method (GAN, Diffusion, etc).
\end{itemize}

% =====================================================================
\section{Dataset Description}
\label{sec:dataset}

The goal was to train, evaluate, and test generalization across three different datasets [listed below]:

\subsection{Dataset \#1: GAN-Based Dataset (FF++)}
Derived from FaceForensics++ \cite{rossler2019faceforensics}, which
contains videos manipulated by four GAN-based methods: Deepfakes,
Face2Face, FaceSwap, and NeuralTextures. 16 frames were extracted from each real video, while 4 frames were extracted from each fake video from each GAN-forgery method. Extracted frames totaled approximately
16\,000 images were split with a 70/15/15 train/test/val.

\subsection{Dataset \#2: Diffusion-Based Dataset}
This dataset was custom built from several different sources. The intended purpose of this dataset was provide us insights into how each architecture would generalize against this dataset when trained on Dataset \#1, as well as generalization against Dataset \#1, when the models were trained on Dataset \#2. This dataset conists of 51,150 images, perfectly balanced between real and fake. The dataset was split with a 70/15/15 train/test/val. Fake images were generated from techniques including Freedom\_T, HPS, MidJourney, SDXL, Freedom\_I, DreamBooth, SDXL\_Refine, DiffFace, CoDiff, Imagic, cycle\_diff, text2img, and inpainting.
Custom built from the below sources:
\begin{itemize}
  \item \textbf{DiFF} \cite{he2021diff}: A collection of
        diffusion-generated face images covering multiple
        sub-categories (HPS, SDXL, DreamBooth, SDXL-Refine, Imagic).
        Only images with a shortest side $\geq 256$\,px are retained.
  \item \textbf{DeepFakeFace} \cite{deepfakeface2023}: Provides
        text-to-image fakes (Stable Diffusion 1.5) and inpainting
        fakes, with real images from IMDB-WIKI celebrities serving as
        the real class.
  \item \textbf{FaceForensics++} \cite{rossler2019faceforensics}: Real
        face images drawn from the genuine video frames of this dataset
        supplement the real class.
  \item \textbf{Celeb-DF-v2} \cite{li2020celeb}: Additional real
        celebrity face images sourced from the real-video splits of
        this benchmark further balance the real class.
\end{itemize}

\subsection{Held-Out Generalization Set}
Celeb-DF-v2 \cite{li2020celeb} is withheld entirely from training and
validation for the official evaluations against models trained on dataset\#1 FF++.  It serves as an unseen-domain test to measure
cross-dataset generalisation.

\paragraph{Implementation Note.} The final and correct evaluation presented in this work were conducted on models trained on the FF++ dataset. We did include dataset \#2  training and evaluation results to show our efforts but these results should not be considered for grading. For dataset\#2, models reached $\sim$99\% AUC within the first couple of epoch of training and after several iterations of reconstructing the dataset and re-training, we were not getting expected results. I included more information in Section~\ref{sec:conclusion}.

% =====================================================================
\section{Data Processing}
\label{sec:processing}

\subsection{Frame Extraction(from sources with mp4)}
Frames were extracted using OpenCV's 'VideoCapture'.The quantity of frames extracted from each mp4 file depended primarily on how many images were needed for the particular partition or deepfake technique the video was derived from. For example, dataset\#1 was created solely from FaceForensics+ so, in order to have a balanced dataset, we needed to extract less frames from each individual forgery technique, which would add up to the amount we extracted from the real videos. The sampling strategy ensured uniform temporal coverage, without bias — frame indices were computed with \texttt{numpy.linspace(0, total\_frames-1, n\_frames)}.

\subsection{Face Detection and Alignment}
MTCNN \cite{Zhang_2016} was used to detect and align faces, with a 20\, \% margin around each detected face. If MTCNN could not detect a face, the image would be opened with PIL and cropped by \texttt{crop\_size = min(width, height)}, and than crops equal margins from both sides of the longer dimension. This center-crop fallback was a valid option since all sourced datasets already guarantee facial structures to be centered but we decided to still use the face detection algorithm to ensure best results.

\subsection{Resolution Filtering}
Any image whose shortest side is $< 256$\,px was discarded.  This was to ensure the images in our dataset were of the best quality and to prevent the model learning artifacts created for smaller images requiring to be upscaled.

\subsection{Balanced Sampling}
Each deepfake-generation method had the exact amount of images in each dataset and split. This prevented model bias towards any particular deepfake generation method by ensuring that no single method held more weight during training, validation, or evaluation.

\subsection{Train / Validation / Test Split}
Each dataset is split into 70\,\% train, 15\,\% validation, and
15\,\% test, honoring the splits provided by the dataset source (if provided). Furthermore, any image or video from the data sources provided train/test/val split would go into the same split of our custom train/test/val splits. 

\subsection{Identity Leakage}
FF++ and most diffusion sources included identity metadata by either providing dataset splits or including identification numbers in the video name. For Celeb-dfv2, the Youtube-Real partition does not have any identity metadata but the creator guarantees that the same person is not found in more than one provided video. Considering this, the Youtube-Real videos were split 70/15/15 seperately (240/45/45 videos) and then frames were extracted from each split to add into the real partition of Diffusion Dataset \#2.

\subsection{Augmentation}
Applied during training only:
\begin{itemize}
  \item ColorJitter: brightness=0.2, contrast=0.2, saturation=0.1
  \item JPEG Compression: 70 to 95
  \item GaussianBlur: \texttt{kernel\_size=3, p=0.3}
  \item Per-channel normalisation: Computed stats for the respective training dataset
  \item For ViT-b16 pretrained: Resize(224) and ImageNet stats
\end{itemize}
Validation and test sets receive only center-crop and normalization. For the pre-trained ViT, images were all resized to 224x224.

% =====================================================================
\section{Methodology}
\label{sec:methodology}

\subsection{Baseline Model}
\label{sec:baseline}

We used the original Mesonet4 and MesoInception4 architecture by Afchar
\emph{et al.}~\cite{afchar2018mesonet} as our baseline.  MesoNet was
designed specifically for detecting face manipulations in video, and
its compact design made it suitable for rapid
experimentation and comparison. Mesoinception4 is a variant of Mesonet, with the first two convolutional blocks being replaced with inception blocks.

Mesonet4 takes a $256 \times 256 \times 3$ input, passes it through four convolutional blocks, with each convolution followed by batch normalization and max pooling. The convolutional blocks are followed by successive fully connected heads with dropout(0.5) and a sigmoid output.

The Mesoinception4 network 
replaces the first two convolution blocks of MEsonet4 with two \emph{inception-like} blocks (each block uses parallel
$1\times1$, $3\times3$, and $5\times5$ convolutions) followed by two
plain convolutional layers, fully
connected heads with a sigmoid output.

\subsection{CNN Model Architecture}
\label{sec:architecture}

My proposed model augments MesoInception4 with Convolutional Block
Attention Modules (CBAM) \cite{woo2018cbam}.  CBAM applies two
sequential attention sub-modules after a convolutional layer:

\begin{enumerate}
  \item \textbf{Channel Attention}: a squeeze-and-excitation style
        gate that re-weights feature channels by pooling spatial
        information (both average and max) and passing through a shared
        MLP.
  \item \textbf{Spatial Attention}: a gate that re-weights spatial
        locations by pooling channel information and convolving with a
        $7 \times 7$ kernel.
\end{enumerate}

CBAM modules are inserted after each of the two inception blocks and
the remaining two convolutional blocks, giving the network the ability to
progressively focus on semantically informative regions (e.g.\ eye region, mouth,
blending boundaries).The architecture that had the best performance is summarized in Table~\ref{tab:architecture}.

\begin{table}[h]
\centering
\caption{MesoInception4-CBAM Architecture Summary}
\label{tab:architecture}
\begin{tabular}{llc}
\toprule
\textbf{Block} & \textbf{Output Size} & \textbf{CBAM} \\
\midrule
Inception Block 1  & $128 \times 128 \times 11$ & \checkmark \\
Inception Block 2  & $64  \times 64  \times 12$ & \checkmark \\
Conv Block 3       & $32  \times 32  \times 16$ & \checkmark \\
Conv Block 4       & $8   \times 8   \times 16$ & \checkmark \\
Flatten            & 1024                       &            \\
FC (1024$\to$16$\to$1) & 1 (logit output)      &            \\
\bottomrule
\end{tabular}
\end{table}

My partners implemented the following architectures:
\begin{enumerate}
  \item \textbf{Hieu Nguyen}: LVNet and the baseline Mesonet4
  \item \textbf{Tajaddin Gafarov}: MesoXceptionCBAM and MesoXceptionNet
\end{enumerate}
\subsection{Architecture Comparison}

Table~\ref{tab:comparison} compares our best model against partner architectures
and pretrained foundation model baselines fine-tuned on the same FF++ training set.

\begin{table}[h]
\centering
\caption{Architecture Comparison --- FF++ Test AUC}
\label{tab:comparison}
\resizebox{\columnwidth}{!}{%
\begin{tabular}{lcc}
\toprule
\textbf{Model} & \textbf{Params} & \textbf{Pretrained} \\
\midrule
MesoNet4 (Hieu Nguyen)  & \textasciitilde{}27K                  & No  \\
MesoInception4 (Ryan Powers) & \textasciitilde{}27K  & No  \\
MesoInception4-CBAM (Ryan Powers) & \textasciitilde{}28K  & No  \\
MesoXceptionCBAM (Tajaddin Gafarov)   & \textasciitilde{}19K  & No  \\
MesoXceptionNet (Tajaddin Gafarov)    & \textasciitilde{}18K  & No  \\
LvNet (Heiu Nguyen) & \textasciitilde{}63M & No \\
Xception \cite{chollet2017xception}  & 22.9M & Yes \\
ViT-B/16 \cite{dosovitskiy2020vit}    & 86M   & Yes \\
\bottomrule
\end{tabular}%
}
\end{table}

% =====================================================================
\subsection{Training Details and Reproducibility}
\label{sec:training}

All experiments are implemented in PyTorch.  Key hyperparameters for my implementation are
summarised in Table~\ref{tab:hparams}.

\begin{table}[h]
\centering
\caption{Training Hyperparameters}
\label{tab:hparams}
\resizebox{\columnwidth}{!}{%
\begin{tabular}{lc}
\toprule
\textbf{Hyperparameter} & \textbf{Value} \\
\midrule
Optimiser          & AdamW (with default beta1 and beta2) \\
Learning rate      & 5e-3 \\
Learning rate schedule & ReduceLROnPlateau with linear warmup (optional arg)\\
Weight decay (L2)  & $1 \times 10^{-4}$ \\
Batch size         & 64 \\
Epochs             & 50 \\
Loss function      & Binary cross-entropy with logits \\
Input resolution   & $256 \times 256$ \\
\bottomrule
\end{tabular}%
}
\end{table}

Early stopping is applied with patience of 10 epochs, monitoring
validation AUC.  Model checkpoints are saved whenever validation AUC
improves and every 10 epochs. All experiments were run on a single NVIDIA GPU with 16gb VRAM.

% =====================================================================
\section{Performance Improvement Techniques}
\label{sec:improvements}

Performance improvements were found with the following techniques:

\begin{itemize}
  \item \textbf{CBAM attention}: detailed in Section~\ref{sec:architecture}.
  \item \textbf{JPEG compression augmentation}: to simulate social-media compression and improve generalization
  \item \textbf{L2 regularisation (weight decay)}: reduces over-fitting
        on the smaller diffusion dataset.
  \item \textbf{Learning rate schedule}: Performance scheduling via Torch.ReduceLRonPlateau
  \item \textbf{Dropout}: Dropout(0.3-0.5) depending on the other regularization techniques and parameters.
\end{itemize}

\begin{table}[h]
\centering
\caption{Results from finetuning and testing hyperparameter techniques}
\label{tab:ablation}
\begin{tabular}{lcc}
\toprule
\textbf{Model Variant} & \textbf{FF++ AUC} & \textbf{$\Delta$ vs.\ best} \\
\midrule
MesoInception4 (baseline)                  & 0.7621          & $-$0.048 \\
\midrule
MesoInception4-CBAM \textbf{(best)}        & \textbf{0.8100} & ---      \\
\quad No CBAM in Inception block~1         & 0.7406          & $-$0.069 \\
\quad No geometric augmentation            & 0.7005          & $-$0.110 \\
\quad No augmentation                      & 0.5772          & $-$0.233 \\
\quad No dropout                           & 0.6749          & $-$0.135 \\
\quad No BatchNorm                         & 0.5004          & $-$0.310 \\
\quad No L2 regularisation                 & 0.6663          & $-$0.144 \\
\quad 100 epochs (over-trained)            & 0.6379          & $-$0.172 \\
\bottomrule
\end{tabular}
\end{table}

\textbf{BatchNorm and augmentation are load-bearing.}
Removing BatchNorm causes complete training failure (AUC~0.50, predicted all
real). Removing all augmentation drops performance near chance (AUC~0.58).
Geometric augmentation gave a $-$0.11 AUC penalty when removed from specific experiments with attention mechanisms,
suggesting that spatial variability during training aids generalization.
% =====================================================================
\section{Performance Evaluation}
\label{sec:evaluation}

\subsection{Metrics}
We report the following metrics for all experiments below. 
\begin{itemize}
  \item \textbf{Accuracy}: overall classification accuracy.
  \item \textbf{AUC-ROC}: area under the receiver operating
        characteristic curve — the primary metric, as it is threshold
        independent.
  \item \textbf{Precision / Recall / F1}: to expose class-specific
        performance.
\end{itemize}

\subsection{Training and Evaluation Information}
This table includes the results from all successful experiments from each team member. Please note the following:
\begin{itemize}
\item Each entry that is trained on FF++ should be considered for evaluation purposes, as the FF++ dataset is verifably reliable and our concerns regarding Dataset \#2 (Diffusion) explained in \autoref{sec:error}.
\item We included the experiments on Dataset\#2 (Diffusion) to show our efforts towards training and evaluating on different deepfake techniques to see if our architectures and techniques saw improvements with learning artifacts from the different forgery methods and generalization.
\end{itemize}

\begin{table*}[t]
\centering
\caption{In-distribution results: models evaluated on the same dataset they were trained on, sorted by AUC.}
\label{tab:ablation_indist}
\begin{tabular}{llccccc}
\toprule
\textbf{Model} & \textbf{Dataset} & \textbf{AUC} & \textbf{Acc.} & \textbf{Prec.} & \textbf{Rec.} & \textbf{F1} \\
\midrule
LVNet                  & Diffusion & 0.9991 & 0.9842 & 0.9770 & 0.9917 & 0.9843 \\
MesoInception4-CBAM    & Diffusion & 0.9935 & 0.9527 & 0.9702 & 0.9341 & 0.9518 \\
MesoInception4         & Diffusion & 0.9926 & 0.9470 & 0.9494 & 0.9442 & 0.9468 \\
Xception (pretrained)  & FF++      & 0.9747 & 0.9179 & 0.9457 & 0.8866 & 0.9152 \\
ResNet-18              & FF++      & 0.9578 & 0.8861 & 0.8643 & 0.9160 & 0.8894 \\
ViT-B/16               & FF++      & 0.9428 & 0.8688 & 0.8918 & 0.8393 & 0.8648 \\
LVNet                  & FF++      & 0.8960 & 0.8217 & 0.8310 & 0.8105 & 0.8206 \\
MesoInception4-CBAM    & FF++      & 0.8100 & 0.7183 & 0.7043 & 0.7527 & 0.7277 \\
MesoInception4         & FF++      & 0.7621 & 0.6933 & 0.7104 & 0.6527 & 0.6803 \\
MesoXceptionNet        & FF++      & 0.7195 & 0.6473 & 0.6315 & 0.7071 & 0.6672 \\
MesoXceptionCBAM       & FF++      & 0.6322 & 0.5941 & 0.5813 & 0.6732 & 0.6239 \\
MesoNet4               & Diffusion & 0.6199 & 0.5171 & 0.5101 & 0.8725 & 0.6438 \\
MesoNet4               & FF++      & 0.6063 & 0.6020 & 0.5840 & 0.6300 & 0.6061 \\
\bottomrule
\end{tabular}
\end{table*}

\begin{table*}[t]
\centering
\caption{Generalization results: models evaluated on a held-out dataset different from their training set, sorted by AUC.}
\label{tab:ablation_gen}
\begin{tabular}{llccccc}
\toprule
\textbf{Model} & \textbf{Train} $\to$ \textbf{Eval} & \textbf{AUC} & \textbf{Acc.} & \textbf{Prec.} & \textbf{Rec.} & \textbf{F1} \\
\midrule
ViT-B/16               & FF++ $\to$ Celeb-DF v2     & 0.6438 & 0.4044 & 0.8369 & 0.1151 & 0.2023 \\
ViT-B/16               & FF++ $\to$ Diffusion       & 0.5850 & 0.5422 & 0.6032 & 0.2468 & 0.3503 \\
MesoInception4         & FF++ $\to$ Celeb-DF v2     & 0.5677 & 0.5208 & 0.7107 & 0.4552 & 0.5550 \\
Xception (pretrained)  & FF++ $\to$ Celeb-DF v2     & 0.5663 & 0.4201 & 0.7361 & 0.1815 & 0.2912 \\
Xception (pretrained)  & FF++ $\to$ Diffusion       & 0.5017 & 0.4807 & 0.4666 & 0.2695 & 0.3416 \\
MesoInception4-CBAM    & Diffusion $\to$ FF++       & 0.4945 & 0.4964 & 0.2500 & 0.0036 & 0.0070 \\
MesoInception4-CBAM    & FF++ $\to$ Celeb-DF v2     & 0.4813 & 0.4886 & 0.6454 & 0.4902 & 0.5572 \\
MesoInception4         & Diffusion $\to$ FF++       & 0.4792 & 0.4987 & 0.4348 & 0.0089 & 0.0175 \\
MesoInception4-CBAM    & FF++ $\to$ Diffusion       & 0.4731 & 0.4721 & 0.4466 & 0.2330 & 0.3062 \\
MesoInception4         & FF++ $\to$ Diffusion       & 0.4556 & 0.4970 & 0.4937 & 0.2333 & 0.3168 \\
LVNet                 & FF++ $\to$ Celeb-DF v2      & 0.570 & 0.49    & --     & --     & 0.485  \\
\bottomrule
\end{tabular}
\end{table*}
% =====================================================================
\section{Comparison of Approaches}
\label{sec:comparison}

\subsection{Generalization}
The generalization problem is universal, including the fine-tuned foundational models. \textbf{Every model fails to generalize on different deepfake methods}. Every model trained on GAN-era deepfakes and tested on diffusion deepfakes has around a 0.45-0.59 AUC, which is effectively guessing. Similar performance was observed for cross-GAN generalization (Celeb-dfv2), with the best results is the pre-trained ViT at 0.64AUC. MesoInception4 trained on Diffusion and tested on FF++ collapses to 0.479 AUC with a recall of 0.009 — it's predicting nearly everything as real. These results show that the models are learning artifacts that are created from specific deepfake generation methods, and the artifacts are not shared between methods or generations of deepfakes.

\subsection{CBAM Impact Analysis}
The impact of CBAM is inconsistent. For my implementation, mesoinception4-cbam, it improved the models performance against regular mesoinception4 on in-domain deepfakes (+~0.048 AUC and F1) but it decreased performance when the FF++ trained model was evaluated against Celebdf-v2. For MesoXception-net, the separable convolutions does not appear to benefit the model, likely because the attention mechanism may be adding too much noise throughout the model or overfitting. The results show that the benefits of CBAM is architecture and data dependent but generally suggests that the network's attention
to spatially localised artifacts is beneficial.

\subsection{Xception dominates in-domain}
Fine-tuned Xception achieves 0.9747 AUC, which shows the advantage of
ImageNet pretraining and larger models. However, this comes at roughly
800$\times$ the parameter cost (\textasciitilde{}28K for mesoinception4 vs.\ 22.9M). The performance of this model in our experiments proves why it has been historically considered the best deepfake detection technique.

\subsection{MesoXceptionNet and MesoXceptionCBAM}
These architectures are the weakest on FF++, with MesoXceptionNet at 0.720 AUC and MesoXceptionCBAM at 0.632. These results hint at maybe an architectural compatibility or training issue. The seperable convolutions sacrifice capacity fort efficiency and maybe the early convolutions are not extracting powerful enough features maps for CBAM to spatially or channelly concentrate on.

\subsection{LvNet}
LVnet appears to perform the best on in-domain evaluations (when compared to the CNN-basedx models that were not pretrained), with a 0.896 AUC for FF++. The performance of this model shows the benefits of a two-stream architecture designed for deepfake detection.

\subsection{Mesonet4 Baseline}
The baseline architecture ended up being the weakest model against FF++, at 0.606AUC. These results show that the legacy architecture is not robust against data compression and augmentation. This also shows that the addition of inception blocks provide critical enhancements in performance, when comparing to the performance of mesoinception4.

% =====================================================================
\section{Foundation Model Comparison}
\label{sec:foundation}

The foundation models include ViT-B16, Resnet18, and Xception. Each model was downloaded via timm as pretrained, and fine-tuned using dataset~\#1 (FF++).

\subsection{ViT-B16}
ViT generalized best out of the pre-trained models, suggesting that the global self-attention captures more facial inconsistencies than the CNN kernels. ViT achieved in-domain 0.9428 AUC and 0.6438 AUC when evaluated against celeb-dfv2. This suggests that even a powerful model with 86M parameters struggles with identifying artifacts across different deepfake generation methods.

\subsection{ResNet-18}
Powerful in-domain, which proves that ImageNet pretraining gives a significant advantage over random initialization.

\subsection{Xception}
Fine-tuned Xception achieves 0.9747 AUC, which shows the advantage of
ImageNet pretraining and larger models. However, this comes at roughly
800$\times$ the parameter cost (\textasciitilde{}28K for mesoinception4 vs.\ 22.9M). The performance of this model in our experiments proves why it has been historically considered the best deepfake detection technique.
 
% =====================================================================
\section{Error and Failure Analysis}
\label{sec:error}

\subsection{Confusion Analysis}
On the FF++ test set, MesoInception4-CBAM achieves 0.7183 accuracy (AUC 0.8100) with recall 0.7527 and precision 0.7043.  The primary failure mode is an elevated
false-negative rate: the model misses fake images more often than it misclassifies real ones (recall $<$ precision gap is consistent across training runs). Per-method, FaceSwap and NeuralTextures are the hardest GAN
methods to detect; Deepfakes-style identity-swap artifacts are comparatively easier.

Table~\ref{tab:permethod_ff} The below table compares per-class accuracy on the FF++ test set
across all four architectures I implemented, trained on FF++.
FaceSwap and NeuralTextures are the hardest manipulation methods for every
model, while Deepfakes and Face2Face are easier. Xception performed best with an overall 0.918 ACC.

\begin{table*}[t]
\centering
\caption{Per-Class Test Accuracy on FF++ (all models trained on FF++).
            Best result per row in \textbf{bold}.}
\label{tab:permethod_ff}
\begin{tabular}{lcccc}
\toprule
\textbf{Class} & \textbf{MesoInception4} & \textbf{MesoInception4-CBAM} & \textbf{ViT-B/16} & \textbf{Xception} \\
\midrule
Deepfakes       & 0.843 & 0.864 & 0.886 & \textbf{0.893} \\
Face2Face       & 0.757 & 0.779 & 0.879 & \textbf{0.896} \\
FaceSwap        & 0.446 & 0.714 & 0.882 & \textbf{0.946} \\
NeuralTextures  & 0.564 & 0.654 & 0.711 & \textbf{0.811} \\
Real            & 0.734 & 0.684 & 0.898 & \textbf{0.949} \\
\midrule
\textbf{Overall} & 0.693 & 0.718 & 0.869 & \textbf{0.918} \\
\bottomrule
\end{tabular}
\end{table*}

\subsection{Generalization Error}
Every model that fails to generalise does so in the same way: it calls almost everything real.
Precision stays relatively high --- when the model does predict fake, it is usually correct ---
but recall collapses, meaning the vast majority of fakes are silently misclassified as real.
Table~\ref{tab:gen_collapse} shows this breakdown across all four architectures evaluated on
Celeb-DF-v2.

\begin{table}[h]
\centering
\caption{Real vs.\ fake accuracy and recall on Celeb-DF-v2 (all models trained on FF++).
         Higher model capacity correlates with more severe fake-recall collapse.}
\label{tab:gen_collapse}
\begin{tabular}{lccc}
\toprule
\textbf{Model} & \textbf{Real Acc.} & \textbf{Fake Acc.} & \textbf{Recall} \\
\midrule
ViT-B/16            & \textbf{95.5\%} & 11.5\% & 0.115 \\
Xception            & 85.6\%          & 18.2\% & 0.182 \\
MesoInception4      & 64.0\%          & 45.5\% & 0.455 \\
MesoInception4-CBAM & 48.6\%          & \textbf{49.0\%} & \textbf{0.490} \\
\bottomrule
\end{tabular}
\end{table}

The pattern is counter-intuitive: the larger and stronger the model, the \emph{more extreme}
the collapse. ViT and Xception are nearly perfect on real images yet almost completely
blind to unseen fakes. The compact MesoInception4 models, while weaker in-domain, are at
least closer to random on out-of-distribution data. This suggests that higher-capacity
models over-fit more tightly to the GAN-specific artefacts present in FF++ training data,
leaving them with no signal when those artefacts are absent.

Generalisation to Celeb-DF-v2 breaks down entirely for both MesoInception variants, as
visualised in Figure~\ref{fig:confusion_celeb}. The confusion matrices confirm the
asymmetric failure: a large off-diagonal mass of fakes predicted as real, with comparatively
few real images mislabelled.
\begin{figure*}[t]
  \centering
  \begin{subfigure}[t]{0.48\linewidth}
    \centering
    \includegraphics[width=\linewidth]{../results/mesoinception4/mesoinception4_trained-FF_eval-Celeb/confusion_matrix.png}
    \caption{MesoInception4 (AUC~0.568, Recall~0.455).
             Predicts real correctly 64\% of the time but misses over half of all fakes.}\label{fig:confusion_celeb_meso}
  \end{subfigure}
  \hfill
  \begin{subfigure}[t]{0.48\linewidth}
    \centering
    \includegraphics[width=\linewidth]{../results/mesoinception4_with_cbam/mesoinception4_cbam_trained-FF_eval-CelebDFv2/confusion_matrix.png}
    \caption{MesoInception4-CBAM (AUC~0.481, Recall~0.490).
             CBAM adds no generalisation benefit; performance degrades to near-random.}\label{fig:confusion_celeb_cbam}
  \end{subfigure}
  \caption{Confusion matrices for models trained on FF++ and evaluated on the held-out
           Celeb-DF-v2 generalisation set. Both models exhibit the same failure mode:
           they call the majority of unseen fakes \emph{real}, confirming that GAN-trained
           detectors over-fit to generation-specific artefacts. CBAM does not improve
           and slightly worsens cross-dataset generalisation (AUC $0.568 \to 0.481$).}\label{fig:confusion_celeb}
\end{figure*}

\textbf{Data scale mismatch.}
ViT was originally pretrained on ImageNet-21k ($>$14M images) and relies on task-specific data to fine-tune its attention patterns
effectively.  The FF++ training set ($\approx$16K images) is insufficient
to adapt 86M parameters without overfitting, regardless of the low learning
rate.  Deepfake detection datasets are orders of magnitude smaller than the
data regimes where ViT excels.

\textbf{Inductive biases matter.}
Convolutional models like MesoInception4-CBAM have built-in spatial locality
and biases well-suited to the task: deepfake
artifacts are local, spatially consistent features (blending boundaries,
texture discontinuities).  ViT's global self-attention has no such prior
and must learn locality from scratch on our limited data. 

% =====================================================================
\section{Team Collaboration}
\label{sec:collaboration}

As a team, we worked together to design the pipeline and the experiments we wanted to perform. Each of us performed different experiments but most of then were related in some way. This helped us learn and understand which architecture or techniques helped detect deepfake artifacts and why.
% =====================================================================
\section{Individual Contributions}
\label{sec:contributions}
\subsection{Myself}
\begin{itemize}
  \item Designed and constructed Dataset~\#1 and Dataset~\#2 and associated manifest files.
  \item Took responsibility for troubleshooting diffusion dataset~\#1
  \item Implemented mesoinception4 and mesoinception4 variants.
  \item Fine-tuned and evaluated the foundation models: ViT-B16 and Xception.
\end{itemize}
\subsection{Hieu}
\begin{itemize}
  \item Implemented Mesonet4 and LVnet
  \item Implemented and evaluated Resnet18
  \item Evaluated models for generalization
\end{itemize}
\subsection{Tajaddin}
\begin{itemize}
  \item Implemented and evaluated MesoxceptionNet and MesoXceptionCBAM
  \item Finetuned and evaluated Resnet18
\end{itemize}

% =====================================================================
\section{Conclusion}\label{sec:conclusion}

We presented a study of deepfake detection generalisation across GAN
and diffusion generation families using an attention-enhanced
MesoInception4 architecture.  Our key findings are:

\begin{itemize}
  \item CBAM improves in-domain but not generalization. Although the findings were inconsistent across CBAM additions between partners, CBAM generally appears to hurt generalization by decreasing AUC.
  \item All models fail to generalize and exhibit the same collapse: high real-accuracy, near-zero fake recall on generalization evaluation.
  \item As model size increases, they magnitude over overfitting increases: ViT with 86M parameters achieves 11.5 percent fake accuracy on celeb-dfv2, while mesoinception4-CBAM with 28k parameters achieves 49 percent. These results show that complex models are not always more effective than smaller models.
  \item Generalizing on diffusion-generated deepfakes is harder than cross-GAN generalization
  \item Neural Textures appears to be the hardest method to detect. This is due to the subtle artifacts this method leaves behind, compared to the larger spatial artifacts from the other methods like Deepfakes and Face2face.
  \item MesoInception4-CBAM achieves better accuracy at $\approx$3{,}000$\times$ fewer parameters and trains in under 2 minutes per epoch versus 70 seconds per epoch for ViT on the same hardware.  Fo constrained deployment scenarios, the compact model is the clear practical choice.
\end{itemize}

I learned a lot about  dataset design and all of the different deepfake techniques. With our diffusion dataset2, Every architecture was achieving .95+ AUC: I knew there was a problem but I did not know exactly what. Originally, I thought it had to do with the original diffusion dataset we were using being perfectly balanced but the real partition only having real images from IMDB-WIKI. I thought a possible solution would be adding real images from multiple sources to poentially prevent bias, so I added real images from the FF++ dataset and CelebDFv2 dataset. The new real partition consisted of .35 Celeb-Youtube, .25 celeb-real, .25 from faceforensics videos, and .15 from IMDB-wiki. After re-training and evaluating, each model still achieved extremely high AUC within the first epochs. While the issue is not fully undferstood, I do know that the combining techniques introduced a lot of oppurtunity for bias by combining images and frames extracted from videos, from various sources. I plan on continuing on with combining all of the real and fake images that I have downloaded and creating a detailed manifest so that the training results will show me exactly which deepfake method or source the real or fake image is coming from. 

Future work includes continuing to research the modern day, SOTA architectures being used for deepfake forensics and detection. I plan on finishing the data problem we faced to learn more about data engineering, the different deepfake generation methods, and how the models studied in this project and beyond perform. I plan on experimenting with other datasets that I found throughout troubleshooting our dataset issue, that has a great distribution of deepfakes from all generations, across all modalities, to see how the models we trained during this project peform. Ultimately, I hope to discover techniques that enable a model to learn artifacts from all manipulation methods and I want to find the best model size vs speed vs accuracy trade off.


% =====================================================================
\bibliographystyle{IEEEtran}
\bibliography{references}

\end{document}
